
\documentclass[a4paper,12pt]{letter}
\usepackage[utf8]{inputenc}
\usepackage[english,slovene]{babel}
%\usepackage{blindtext}
%\usepackage{pdflscape}
\usepackage[tikz]{bclogo}
\usepackage{qrcode,marvosym}
\usepackage{pgf,tikz,pgfplots}
\pgfplotsset{compat=1.14}
\usepackage{mathrsfs}
\usetikzlibrary{arrows}
\usepackage{graphicx}
\newcommand{\ds}{\displaystyle}
\usepackage{fancybox}
\usepackage{amsfonts}
\usepackage{amsmath}
%\usepackage{color}
\usepackage{wallpaper}
\usepackage{hyperref}
\usepackage[cp1250]{inputenc}
\usepackage{array}%%%%%%%%%%%%%%%%%%%%%%%%%%%%%%%%%%%%%%%%%%%%
\newcolumntype{M}[1]{>{\centering\arraybackslash}m{#1}}%%%%
\newcolumntype{N}{@{}m{0pt}@{}}%%%%%%%%%%%%%%%%%%%%%%%%%%%%
   \usepackage{fancyhdr,enumerate}

 \newcommand{\ora}{\overrightarrow}
 \renewcommand{\headrulewidth}{3pt}
 \renewcommand{\headrule}{\hbox to\headwidth{%
   \color{gray}\leaders\hrule height \headrulewidth\hfill}}
    \renewcommand{\footrulewidth}{2pt}
    \renewcommand{\footrule}{\hbox to\headwidth{%
  \color{brown}\leaders\hrule height \footrulewidth\hfill}}
\usepackage{gensymb}
\usepackage{amssymb} 

  \usepackage{fancyhdr}
  \textwidth 20 true cm
  \oddsidemargin -2 true cm
  \evensidemargin -2 true cm
  \topmargin -3.5 true cm
  \textheight 27.5 true cm
  \linespread{1.2}
    \renewcommand{\headrulewidth}{0.3pt}
   \renewcommand{\footrulewidth}{1.9pt}
  \definecolor{qqwuqq}{rgb}{1,0,0 }
  \definecolor{qqqqff}{rgb}{0,0,1}
  \definecolor{qqffqq}{rgb}{0,1,0}
  \definecolor{ffqqqq}{rgb}{1,0,0}
\definecolor{zzuxux}{rgb}{0,0,0}%{0.92,0.03,0.03}
\usepackage{mdframed}
\usepackage{multicol}
 \definecolor{bgblue}{RGB}{0,0,253}
\usepackage{xcolor}
\usepackage{eurosym}
%%%%%%%%%%%%%%%%%%%%%%%%%%%%%%%%%%%%%%%%%%%%%%%%%%%%%%%%
\newcounter{tocke}
\setcounter{tocke}{0}
\usepackage{tikz-3dplot}
%%%%%%%%%%%%%%%%%%%%%%%%%%%%%%%%%%%%%%%%%%%%%%%%%%%%%%%%%%
\mdfdefinestyle{theoremstyle}{%
linecolor=brown!40,linewidth=1pt,%
frametitlerule=true,%
frametitlebackgroundcolor=brown!50,
innertopmargin=\topskip,
}

\mdfdefinestyle{theoremstyle1}{%
linecolor=brown,middlelinewidth=1pt,%
frametitlerule=true,%
apptotikzsetting={\tikzset{mdfframetitlebackground/.append style={%
shade,left color=black!40, right color=gray!10},
mdfbackground/.append style={
top color=white!100!white,
bottom color=black!5
},
}}, 
frametitlerulecolor=gray,
frametitlerulewidth=2pt,
innertopmargin=\topskip,
innerbottommargin=\topskip,
}
\mdtheorem[style=theoremstyle1]{naloga}{Naloga}
\mdtheorem[style=theoremstyle]{resitev}{Analiza Naloge}
\mdtheorem[style=theoremstyle]{kriterij}{\v{S}tevilo dose\v{z}enih to\v{c}k na testu}
%%%%%%%%%%%%%%%%%%%%%%%%%%%%%%%%%%%%%%%%%%%%%%%%%%%%%%%%%%
\usepackage[table]{xcolor}
\usepackage{venndiagram}
\usepackage[printwatermark]{xwatermark}
\usepackage{xcolor}
\newwatermark[allpages,color=brown!2,angle=45,scale=5,
xpos=-5,ypos=0]{matej.info}
%%%%%%%%%%%%%%%%%%%%%%%%%%%%%%%%%%%%%%%%%%%%%%%%%%%%%%%%%%
%%%%%%%%%%%%%%%%%%%%%%%%%%%%%%%%%%%%%%%%%%%%%%%%%%%%%%%%%%
\begin{document}
\pagestyle{fancy}
\renewcommand\headrulewidth{0pt}
\lhead{}\chead{}\rhead{}
\rfoot{\vspace*{-2\baselineskip}}


\renewcommand{\vec}{\overrightarrow}

%%%%%%%%%%%%%%%%%%%%%%55\Large

\everymath{\displaystyle}
\begin{minipage}{20cm}
\begin{bclogo}[couleur=brown!40,
  arrondi=0,ombre=true, couleurOmbre=gray!50,
logo =\bcplume,
  noborder=true]
{\textrm{{\Large{Test 1.1.} \Huge{}:
%
%
%%%%%%%%%%%%%%%%%%%%%%%%%%%%%%%%%%%%%%%%%
%%%%%%%%%%%%%%%%%%%%%%%%%%%%%%%%%%%%%%%%%
% TUKAJ VNESI VSEBINO TESTA
\large{Polinomi in racionalna funkcija} \hfill $G-3$}}}
%%%%%%%%%%%%%%%%%%%%%%%%%%%%%
\end{bclogo}
\end{minipage}
\begin{minipage}{20cm}
\begin{bclogo}[couleur=brown!40,ombre=true, couleurOmbre=gray!50,
  arrondi=0,
logo =\bcplume,
  noborder=true]
 \setlength{\parskip}{}{}
{\textsc{{\large{Ime in priimek:}
%%%%%%%%%%%%%%%%%%%%%%%%%%%%%%%%%%%%%%%%%
% TUKAJ VNESI VSEBINO TESTA
\vspace{0.2cm} \rule{15cm}{0.2mm}\hfill }}}
%%%%%%%%%%%%%%%%%%%%%%%%%%%%%%%%%%%%%%%%%
\end{bclogo}
\end{minipage}


\def\tocke2{3+3}
\begin{naloga}[\hfill \quad $\tocke2$
\qquad \rightsquigarrow \vert a.\qquad|] \addtocounter{tocke}{\tocke2}

Izračunaj s Hormerjevim algoritmom:

a) vrednost polinoma $p(x)= \dfrac{1}{2}x^3-x^2+7x-\dfrac{3}{2}$ v točki $3$,

b) količnik pri deljenju polinoma $p(x)=x^4-2x^3+3x^2-5$ s polinomom $x-2$.
\end{naloga}





\newpage
\def\tocke18{6}
\begin{naloga}[\hfill \quad $\tocke18$
\qquad \rightsquigarrow \vert a.\qquad\vert b.\qquad|
c.\qquad|] \addtocounter{tocke}{\tocke18}

Ali je polinom $p(x)=x^3-4x^2-x+4$ deljiv s polinomom $x^2-5x+4?$
Pokaži z računom.
\end{naloga}

\vspace{10cm}
\def\tocke6{5}
\begin{naloga}[\hfill \quad $\tocke6$
\qquad \rightsquigarrow \vert a.\qquad| b.\qquad| c.\qquad| ] \addtocounter{tocke}{\tocke6}
Določi ničle polinoma $p(x)=x^3-2x^2-5x+6$ s Hornerjem, tako da najprej zapišeš vse kandidate za celoštevilske ničle.
\end{naloga}

\newpage
\def\tocke3{4+3}
\begin{naloga}[\hfill \quad $\tocke3$
\qquad \rightsquigarrow \vert a.\qquad\vert b.\qquad|
c.\qquad|d.\qquad|] \addtocounter{tocke}{\tocke3}
Nariši graf polinoma $p(x)=2(x+1)^2(x-2)$, tako da prej določiš ničle in njihovo stopnjo ter začetno vrednost.

Koliko je $p(-2)?$

\end{naloga}

\vspace{10cm}

\def\tocke1{4}
\begin{naloga}[\hfill \quad $\tocke1$
\qquad \rightsquigarrow \vert a.\qquad\vert b.\qquad|
c.\qquad|d.\qquad|] \addtocounter{tocke}{\tocke1}
Reši enačbo: $x^3-3x=2x^2$
\end{naloga}
\newpage
\def\tocke6{8}
\begin{naloga}[\hfill \quad $\tocke6$
\qquad \rightsquigarrow \vert a.\qquad| b.\qquad| c.\qquad| ] \addtocounter{tocke}{\tocke6}
Določi ničlo, pol in asimptoto funkcije $f(x)=\dfrac{2x-6}{x+1}$ in jo nariši.
\end{naloga}


\vfill
\mdtheorem[style=theoremstyle1]{kriterij}{Kriterij ocenjevanja}
\begin{kriterij*}[\hfill \v{s}tevilo vseh to\v{c}k na testu: $\arabic{tocke}$]
 $$
 \begin{array}{||c|c|c|c|c|c||
>{}c|c|}
 \hline\hline
 \textup{ocena}& 1&2&3&4&5&\textrm{uspe\v{s}nost v } \%& \textrm{\bf OCENA}\\
  \hline\hline
 \% & [0,45) &[45,60)&[60,75
 ) &[75,90)    &    [90,100]& \mbox{\phantom{ \huge{$\frac{5}{445}$
 15}}
 } 
 & \\
 \hline
 \end{array} \hspace*{0.5cm}\vspace*{-0.3cm} \qrcode[hyperlink,height=0.8in]{http://www.matej.info}
$$
\vspace*{0.1cm}
\end{kriterij*}

\newpage


%% RE\v{S}ITVE
%%%%%%%%%%%%%%%%%%%%%%%%%%%%%%%%%%%%%%%%%%%%%%%%%%%%%%%%%%
\newpage
\begin{minipage}{20cm}
\begin{bclogo}[couleur=brown!40,
  arrondi=0,ombre=true, couleurOmbre=gray!50,
logo =\bcplume,
  noborder=true]
{\textrm{{\Large RE\v{S}ITVE -- Test 1.1.:}
\large{ Polinomi in racionalna funkcija} \hfill $G-3$}}
\end{bclogo}
\end{minipage}
 
\begin{resitev*}[Naloga 1 -- Hornerjev algoritem]
 
\textbf{a)} $p(x)=\tfrac{1}{2}x^3-x^2+7x-\tfrac{3}{2}$, \quad i\v{s}\v{c}emo $p(3)$.
 
Koeficienti: $\tfrac{1}{2},\; -1,\; 7,\; -\tfrac{3}{2}$
 
\[
\begin{array}{r|rrrr}
   & \tfrac{1}{2} & -1 & 7 & -\tfrac{3}{2} \\
 3 &              & \tfrac{3}{2} & \tfrac{3}{2} & \tfrac{51}{2} \\
\hline
   & \tfrac{1}{2} & \tfrac{1}{2} & \tfrac{17}{2} & 24
\end{array}
\]
 
$$\boxed{p(3) = 24}$$
 
\textbf{b)} $p(x)=x^4-2x^3+3x^2+0x-5$, \quad delimo z $(x-2)$.
 
\[
\begin{array}{r|rrrrr}
   & 1 & -2 & 3 & 0 & -5 \\
 2 &   &  2 & 0 & 6 & 12 \\
\hline
   & 1 &  0 & 3 & 6 &  7
\end{array}
\]
 
$$p(x) = (x-2)(x^3+3x+6)+7$$
 
Koli\v{c}nik: $\boxed{x^3+3x+6}$, \quad ostanek: $7$.
 
\end{resitev*}
 
\begin{resitev*}[Naloga 2 -- Deljivost polinomov]
 
Ali je $p(x)=x^3-4x^2-x+4$ deljiv z $x^2-5x+4$?
 
Najprej faktoriziramo delitelja: $x^2-5x+4 = (x-1)(x-4)$.
 
Po izreku o ni\v{c}lah je $p(x)$ deljiv z $x^2-5x+4$, \v{c}e velja $p(1)=0$ in $p(4)=0$.
 
Preverimo s Hornerjem:
 
$p(1)$: koeficienti $1,\,-4,\,-1,\,4$
\[
\begin{array}{r|rrrr}
   & 1 & -4 & -1 & 4 \\
 1 &   &  1 & -3 & -4 \\
\hline
   & 1 & -3 & -4 &  0
\end{array}
\]
$p(1)=0$ \checkmark
 
$p(4)$ iz koli\v{c}nika $x^2-3x-4$: koeficienti $1,\,-3,\,-4$
\[
\begin{array}{r|rrr}
   & 1 & -3 & -4 \\
 4 &   &  4 &  4 \\
\hline
   & 1 &  1 &  0
\end{array}
\]
$p(4)=0$ \checkmark
 
Ker sta obe ni\v{c}li, je $p(x)$ deljiv z $x^2-5x+4$.
$$\boxed{p(x) = (x^2-5x+4)(x+1) = (x-1)(x-4)(x+1)}$$
 
\end{resitev*}
 
\begin{resitev*}[Naloga 3 -- Ni\v{c}le s Hornerjem]
 
$p(x)=x^3-2x^2-5x+6$
 
Kandidati za celo\v{s}tevilske ni\v{c}le (delitelji prostega \v{c}lena $6$):
$$\pm1,\;\pm2,\;\pm3,\;\pm6$$
 
Preizkus $x=1$:
\[
\begin{array}{r|rrrr}
   & 1 & -2 & -5 & 6 \\
 1 &   &  1 & -1 & -6 \\
\hline
   & 1 & -1 & -6 &  0
\end{array}
\]
$p(1)=0$ \checkmark
 
Koli\v{c}nik $x^2-x-6 = (x-3)(x+2)$.
 
$$\boxed{p(x)=(x-1)(x-3)(x+2), \quad \text{ni\v{c}le: } x_1=1,\; x_2=3,\; x_3=-2}$$
 
\end{resitev*}
 
\begin{resitev*}[Naloga 4 -- Graf polinoma]
 
$p(x)=2(x+1)^2(x-2)$
 
\textbf{Ni\v{c}le in njihove stopnje:}
\begin{itemize}
  \item $x=-1$ stopnje $2$ (graf se \textit{dotakne} osi $x$, ne pre\v{c}ka)
  \item $x=2$ stopnje $1$ (graf \textit{pre\v{c}ka} os $x$)
\end{itemize}
 
\textbf{Za\v{c}etna vrednost:} $p(0)=2\cdot1^2\cdot(-2)=\boxed{-4}$
 
\textbf{Vodilni koeficient:} $2>0$, stopnja $3$ $\Rightarrow$ graf gre iz $-\infty$ (levo) v $+\infty$ (desno).
 
\textbf{Vrednost $p(-2)$:}
$$p(-2)=2(-2+1)^2(-2-2)=2\cdot(-1)^2\cdot(-4)=2\cdot1\cdot(-4)=\boxed{-8}$$
 
\begin{center}
\begin{tikzpicture}
\begin{axis}[
  width=10cm, height=7cm,
  axis lines=middle,
  xlabel={$x$}, ylabel={$y$},
  xmin=-3, xmax=3.5,
  ymin=-10, ymax=10,
  xtick={-2,-1,0,1,2,3},
  ytick={-8,-4,0,4,8},
  grid=both, grid style={line width=0.3pt, draw=gray!30},
  tick label style={font=\small},
]
\addplot[blue, thick, domain=-2.5:3, samples=100]
  {2*(x+1)^2*(x-2)};
\addplot[only marks, mark=*, mark size=2pt, color=red]
  coordinates {(-1,0) (2,0) (0,-4) (-2,-8)};
\node[above right, font=\small] at (axis cs:-1,0) {$(-1,0)$};
\node[above right, font=\small] at (axis cs:2,0) {$(2,0)$};
\node[right, font=\small] at (axis cs:0,-4) {$(0,-4)$};
\node[left, font=\small] at (axis cs:-2,-8) {$(-2,-8)$};
\end{axis}
\end{tikzpicture}
\end{center}
 
\end{resitev*}
 
\begin{resitev*}[Naloga 5 -- Re\v{s}evanje ena\v{c}be]
 
$$x^3-3x=2x^2$$
$$x^3-2x^2-3x=0$$
$$x(x^2-2x-3)=0$$
$$x(x-3)(x+1)=0$$
 
$$\boxed{x_1=0,\quad x_2=3,\quad x_3=-1}$$
 
\end{resitev*}
 
\begin{resitev*}[Naloga 6 -- Racionalna funkcija]
 
$$f(x)=\frac{2x-6}{x+1}$$
 
\textbf{Ni\v{c}la} (tel.~$= 0$): $2x-6=0 \;\Rightarrow\; x=3$. \quad Ni\v{c}la: $\boxed{x=3}$
 
\textbf{Pol} (im.~$= 0$): $x+1=0 \;\Rightarrow\; x=-1$. \quad Pol: $\boxed{x=-1}$
 
\textbf{Vodoravna asimptota} (primerjamo stopnji -- enaki stopnji):
$$\lim_{x\to\pm\infty} f(x) = \frac{2}{1} = 2 \quad \Rightarrow \quad \boxed{y=2}$$
 
Lahko zapi\v{s}emo tudi: $f(x)=\dfrac{2(x+1)-8}{x+1}=2-\dfrac{8}{x+1}$, kar potrdi $y=2$ in premik standardne hiperbole.
 
\begin{center}
\begin{tikzpicture}
\begin{axis}[
  width=9cm, height=8cm,
  axis lines=middle,
  xlabel={$x$}, ylabel={$y$},
  xmin=-6, xmax=7,
  ymin=-6, ymax=8,
  xtick={-5,-4,-3,-2,-1,0,1,2,3,4,5,6},
  ytick={-4,-2,0,2,4,6},
  grid=both, grid style={line width=0.3pt, draw=gray!30},
  tick label style={font=\small},
]
% veja x < -1
\addplot[blue, thick, domain=-6:-1.2, samples=80]
  {(2*x-6)/(x+1)};
% veja x > -1
\addplot[blue, thick, domain=-0.8:6.5, samples=80]
  {(2*x-6)/(x+1)};
% pol (navpi\v{c}na asimptota)
\addplot[red, dashed, thick] coordinates {(-1,-6) (-1,8)};
% vodoravna asimptota
\addplot[brown, dashed, thick] coordinates {(-6,2) (7,2)};
% oznake
\addplot[only marks, mark=*, mark size=2pt, color=red]
  coordinates {(3,0) (0,-6)};
\node[above right, font=\small] at (axis cs:3,0) {$(3,0)$};
\node[right, font=\small] at (axis cs:0,-6) {$(0,-6)$};
\node[above left, font=\small, red] at (axis cs:-1,7) {$x=-1$};
\node[right, font=\small, brown] at (axis cs:5,2.4) {$y=2$};
\end{axis}
\end{tikzpicture}
\end{center}
 
\end{resitev*}

\end{document}